\documentclass[11pt,a4paper]{article}
\usepackage[T1]{fontenc}\usepackage[utf8]{inputenc}\usepackage{lmodern}
\usepackage[a4paper,margin=2.2cm]{geometry}
\usepackage{amsmath,amssymb}\usepackage{enumitem}\usepackage{booktabs,longtable,array}
\usepackage[hidelinks]{hyperref}
\setlength{\parindent}{0pt}\setlength{\parskip}{0.4em}
\title{\textbf{Big Bang transition front (OP3): opening bridge from the
OP-CONST/OP-grad chain}\\\large Recon v1 --- proposal-only}
\author{Per reviewer directive after F7 acceptance (Variant B)}
\date{12 June 2026}
\begin{document}\maketitle

\section*{0. Scope and firewall}
This document opens the Big Bang transition front. It does \emph{not}
attempt the ordering mechanism (OP3.1, recorded Open with interface
requirements); it does three bounded things: (i)~maps the recorded state of
OP3; (ii)~imports three new structures from the just-closed
OP-CONST/G1/G2/F7 chain into the OP3 sub-problems where they bind;
(iii)~selects the first attackable work items. All statements are
conditional and respect the recorded interface requirements (no
emergent-time or causal structure presupposed before ordering; selection as
Euclidean measure concentration; no energy statements at the pre-emergent
level). Proposal-only; no preprint edits.

\section{Recorded state of OP3 (anchor map)}
The record already provides a sharpened programme:
\textbf{OP3.1} ordering mechanism, decomposed into
\emph{a}~(bare-to-ordered gap; explicitly recorded as coinciding with
completing OP-grad), \emph{b}~(EFT-data derivation: $Z_u$, $\mathcal C_n$,
$W(u)$, hence $F,\omega,U$ of the closure-function map; subsumes OP2),
\emph{c}~(selection as $S_0$-weighted Euclidean measure concentration);
\textbf{OP3.2} branch orientation, partially clarified by the time-reversal
covariance lemma (Level~A), the no-bounce sublemma (Level~A conditional),
and the conditional persistence criterion (Level~B);
\textbf{OP3.3} thermal completion, with a conditional bookkeeping ledger
($\rho_{\rm onset}$, $T_{\rm th}$, $w_{\rm pre}$) and a late-stage
Kibble--Zurek scaffold, sharply distinguished from geometric branch onset;
\textbf{OP3.4} relics (open observational programme).
The Big-Bang singularity is recorded as a branch boundary, with breakdown
indicators (i)~$\rho\to0$, (ii)~$T_{\rm loc}\to T_c$,
(iii)~$\ell_{\rm cond}\sim m_\sigma^{-1}$ reached by curvature/gradient
invariants. The front's task is to populate this recorded skeleton, not to
re-derive it.

\section{Bridge 1: the F7 constraint is a recorded filter on OP3.1b}
OP3.1b asks for the ordered-branch EFT data --- including the amplitude
potential and kinetic normalisation --- derived from bare
$(\mu^2,\lambda,\alpha,\beta)$. The F7 recon decomposed exactly this object
on the gap side: $m_\sigma^2(u)=V''_{\rm eff}(u)/Z_\sigma(u)$, with the
recorded matching network fixing both the value at $u_0$
($m_\sigma\simeq\sqrt{\hat a_{\rm eff}/16\pi^2}\,u_0/\phi_0$) and, under the
recorded $c_G=-1$, the slope
\[
\frac{\partial}{\partial\ln u}\ln\frac{V''_{\rm eff}}{Z_\sigma}\;=\;1
\quad\text{at }u=u_0.
\]
\textbf{Import:} any OP3.1b candidate derivation of
$(Z_u,\mathcal C_n,W(u))$ now faces a quantitative necessary condition
inherited from the constants chain --- it must reproduce this value-and-slope
pair (or explicitly land on the $c_G=-2$ branch and accept the
closure-map renormalisation classified in F7). The transition front and the
OP-grad core thus converge on \emph{one} object, as the record itself
anticipates (OP3.1a $\equiv$ OP-grad), but now with a differential filter
attached rather than only an existence requirement.

\section{Bridge 2: a breakdown-surface discriminator from the F7 branches}
The recorded breakdown indicators include (i)~$\rho\to0$ and
(iii)~$\ell_{\rm cond}\sim m_\sigma^{-1}$. The F7 sector-of-origin branches
give these indicators different relative behaviour as the ordering
threshold is approached ($u\to$ small):
\begin{description}[leftmargin=2.0em,itemsep=2pt]
\item[Gradient-native ($m_\sigma^2=c\,u$, $c_G=-1$):] the UV threshold
collapses together with the order parameter, $\ell_{\rm cond}\propto
u^{-1/2}\to\infty$; indicators (i) and (iii) describe \emph{one merged
breakdown surface} --- the EFT cutoff and the ordering threshold are the
same physical boundary.
\item[Potential-imported (pinned $m_\sigma$, $c_G=-2$):] $\ell_{\rm cond}$
stays fixed while the amplitude vanishes; indicators (i) and (iii) define
\emph{split surfaces}, and an intermediate regime exists where the
amplitude is small but the cutoff is still high.
\end{description}
\textbf{Import:} the $c_G$ diagnostic of F7 acquires a transition-side
face: the structure of the branch boundary (merged vs split breakdown
surfaces) discriminates the same two branches. This is a kernel-level,
conditional classification --- no transition dynamics is solved --- but it
hands OP3 a structural question whose answer simultaneously settles F7.

\section{Bridge 3: charge-selective routing of thermal completion}
The G2/portal architecture constrains OP3.3 after branch onset (all
statements here live strictly on the post-onset side, per the recorded
interface requirements). The drifting/relaxing amplitude mode carries no
recorded holonomy or charge label and, within the recorded induced-gauge
architecture, has no tree-level $\sigma F^2$ portal; loop-induced portals
reduce to the charged-sector selection rule. \textbf{Import (conditional
necessary-sequence constraint for the OP3.3 ledger):} the condensate
relaxation energy cannot be dumped directly into gauge radiation through a
neutral-mode portal; thermalisation must route through the charged matter
chain, whose masses anchor to $m_f=y_fv_2$ with
$v_2=\phi_0e^{-\mathcal I_{\rm EW}}$. Hence any thermal-completion
mechanism inherits an event-ordering requirement:
\emph{gradient ordering $\to$ metric/$\Phi$-sector matching (so that the
$v_2$ chain exists) $\to$ charged-spectrum population $\to$ entropy
production and radiation domination.}
This adds a structural line to the recorded bookkeeping ledger --- a
constraint on mechanisms, with no mechanism implied --- and it makes the
charge-based selectivity of the constants chain do double duty: the same
architecture that protects $\alpha,\mu$ at late times channels the energy
flow at early times.

\section{Work plan and verdict}
\begin{enumerate}[leftmargin=1.9em,itemsep=2pt,label=\textbf{B\arabic*.}]
\item (This document.) Anchor map + three bridges imported from the
OP-CONST/G1/G2/F7 chain into OP3.1b, the branch-boundary structure, and the
OP3.3 ledger. Nothing here advances the OP3.1 mechanism itself.
\item Candidate next recon: formalise the merged-vs-split breakdown-surface
statement of Bridge~2 as a conditional sublemma at the level of the
recorded breakdown criterion (kernel-conditional; would also serve as an
internal F7 discriminator).
\item Candidate next recon: draft the Bridge-3 event-ordering line in the
exact format of the recorded thermal-completion ledger, for possible later
D-insertion as one ledger row (only after review).
\item Honesty clause: OP3.1a--c remain Open exactly as recorded; the
interface requirements are respected throughout (no pre-emergent energy
statements; no time-language for the Euclidean selection problem); all
bridges are conditional on the Level~B kernel and the G2/portal
architecture already accepted at recon level.
\end{enumerate}
\end{document}
